\documentclass[11pt,a4paper]{article}

% === 中文支持 ===
\usepackage{ctex}
\usepackage[T1]{fontenc}

% === 页面布局 ===
\usepackage[top=2.5cm, bottom=2.5cm, left=2.5cm, right=2.5cm]{geometry}

% === 表格 ===
\usepackage{booktabs}
\usepackage{longtable}
\usepackage{array}
\usepackage{enumitem}
\usepackage{tabularx}

% === 代码 ===
\usepackage{listings}
\lstset{
  basicstyle=\ttfamily\footnotesize,
  breaklines=true,
  frame=single,
  backgroundcolor=\color{gray!8},
  showstringspaces=false
}

% === 数学符号 ===
\usepackage{amssymb}

% === 颜色 ===
\usepackage{xcolor}
\definecolor{primary}{HTML}{1a1a2e}
\definecolor{accent}{HTML}{3b82f6}
\definecolor{success}{HTML}{22c55e}
\definecolor{danger}{HTML}{ef4444}
\definecolor{warning}{HTML}{f59e0b}
\definecolor{graytext}{HTML}{6b7280}

% === 超链接 ===
\usepackage[hidelinks]{hyperref}
\hypersetup{colorlinks=true, linkcolor=primary, urlcolor=accent}

% === 标题样式 ===
\usepackage{titlesec}
\titleformat{\section}{\Large\bfseries\color{primary}}{}{0em}{}[\titlerule]
\titleformat{\subsection}{\large\bfseries\color{primary}}{}{0em}{}
\titleformat{\subsubsection}{\normalsize\bfseries\color{primary}}{}{0em}{}

% === 页眉页脚 ===
\usepackage{fancyhdr}
\pagestyle{fancy}
\fancyhf{}
\fancyhead[L]{\small\color{graytext} 企业级 AI Agent 应用商店 -- 项目总纲}
\fancyhead[R]{\small\color{graytext} v0.1}
\fancyfoot[C]{\small\color{graytext} \thepage}
\renewcommand{\headrulewidth}{0.4pt}

% === 文档信息 ===
\title{
\vspace{-1em}
{\Huge\bfseries\color{primary} 企业级 AI Agent 应用商店}\\[0.8em]
{\Large\bfseries\color{accent} 项目总纲}\\[0.5em]
{\large V0 先备阶段 -- 智能体分工、技术选型与文档体系}
}
\author{
\textbf{PM AI}：万能产品小助手（OpenClaw）\\[0.3em]
\textbf{审阅}：魏子政
}
\date{2026-05-06 \quad 版本 v0.1}

\begin{document}

\maketitle
\thispagestyle{empty}

\vspace{1em}

\begin{center}
\begin{tabular}{>{\bfseries}p{3cm} p{10cm}}
\toprule
项目名称 & 企业级 AI Agent 应用商店 \\
文档类型 & 项目总纲（V0 先备阶段） \\
当前版本 & v0.1 \\
文档阶段 & V0（先备阶段：智能体分工 + 整体框架 + 文档体系搭建） \\
技术栈 & Vue 3 + Element Plus + FastAPI + PostgreSQL + MinIO + OpenSearch + JWT + Docker Compose \\
PM AI & 万能产品小助手（OpenClaw） \\
开发者 & Cursor Agent \\
\bottomrule
\end{tabular}
\end{center}

\vspace{1.5em}

\begin{abstract}
本文档是企业级 AI Agent 应用商店项目的总纲，面向项目决策者（魏子政）和未来加入的团队成员。内容覆盖四大核心板块：

\begin{enumerate}[nosep]
  \item \textbf{智能体搭配与 Skill 加载机制}：OpenClaw（PM AI）和 Cursor Agent 的分工、各自 Skill/规则如何加载、两者的本质区别
  \item \textbf{技术选型}：每项技术的选用原因、与备选方案的对比、技术栈锁定规则
  \item \textbf{文档体系与生命周期}：tex-pdf 格式规范、命名规则、版本迭代、需求变更处理
  \item \textbf{自闭环机制}：从需求到交付的五个闭环、监控标准、质量保障
\end{enumerate}

当前项目处于 \textbf{V0 先备阶段}，本总纲是 V0 的核心交付物之一。V0 阶段的目标是搭好智能体分工体系和项目框架，V1 阶段开始接入正式业务需求。
\end{abstract}

\newpage
\tableofcontents
\newpage

% ============================================================
%  第一部分：智能体搭配
% ============================================================
\part{智能体搭配与 Skill 加载机制}

% ---------- 第一章 ----------
\section{为什么需要两种智能体}

本项目是一个完整的企业级应用，需要两类截然不同的能力：

\begin{table}[h]
\centering
\begin{tabular}{p{3cm}p{4cm}p{4cm}}
\toprule
\textbf{能力类型} & \textbf{说明} & \textbf{适合的智能体} \\
\midrule
产品管理 & 需求理解、任务拆解、质量审计、进度追踪、文档维护、飞书沟通 & OpenClaw（PM AI） \\
代码开发 & 编写前后端代码、运行终端命令、构建项目、调试修复 & Cursor Agent \\
\bottomrule
\end{tabular}
\end{table}

\textbf{核心分工链路}：魏子政只跟 OpenClaw 对话 $\to$ OpenClaw 拆解需求并调度 Cursor $\to$ Cursor 写代码 $\to$ OpenClaw 审计质量 $\to$ 向魏子政汇报。

\textbf{为什么不能只用一个？}

\begin{itemize}[nosep]
  \item OpenClaw 是\textbf{常驻 AI 助手}：有记忆系统（memory 文件）、有飞书通道、能主动推送、能执行非代码任务
  \item Cursor 是\textbf{代码编辑器内置 AI}：擅长读写代码、运行命令、理解项目上下文，但没有持久记忆和外部通信能力
  \item 两者互补：OpenClaw 管\textbf{脑}（规划/审计/记忆），Cursor 管\textbf{手}（编码/构建/自测）
\end{itemize}

% ---------- 第二章 ----------
\section{两种智能体的 Skill 加载机制}

这是整个协作体系最核心的区别。

\subsection{OpenClaw 的 Skill 加载}

\begin{itemize}[nosep]
  \item \textbf{触发时机}：OpenClaw 启动会话时，系统提示中包含 \texttt{<available\_skills>} 列表
  \item \textbf{加载方式}：收到用户消息后，AI \textbf{自主判断}该加载哪个 skill，用 \texttt{read} 工具读取对应的 \texttt{SKILL.md}
  \item \textbf{存放位置}：
    \begin{itemize}[nosep]
      \item \texttt{\textasciitilde/.openclaw/skills/} -- 全局 skills（所有项目共享）
      \item \texttt{workspacenew1/skills/} -- 工作区 skills（项目专属）
    \end{itemize}
  \item \textbf{文件格式}：普通 Markdown，文件名固定 \texttt{SKILL.md}
\end{itemize}

\subsection{Cursor 的规则加载}

\begin{itemize}[nosep]
  \item \textbf{触发时机}：Cursor 打开项目工作区时自动扫描
  \item \textbf{加载方式}：
    \begin{itemize}[nosep]
      \item \texttt{CLAUDE.md} -- 项目根目录，\textbf{每次对话自动加载}，总入口
      \item \texttt{.cursor/rules/*.mdc} -- MDC 格式，支持两种触发：
        \begin{itemize}[nosep]
          \item \texttt{alwaysApply: true} $\to$ 每次对话都加载
          \item \texttt{globs: ["backend/**"]} $\to$ 只在编辑匹配目录的文件时加载
        \end{itemize}
    \end{itemize}
  \item \textbf{文件格式}：MDC（YAML front matter + Markdown）
\end{itemize}

\subsection{两者的本质区别}

\begin{table}[h]
\centering
\small
\begin{tabular}{p{2.5cm}p{4.5cm}p{4.5cm}}
\toprule
\textbf{维度} & \textbf{OpenClaw Skill} & \textbf{Cursor Rule} \\
\midrule
文件位置 & \texttt{skills/xxx/SKILL.md} & \texttt{.cursor/rules/xxx.mdc} \\
给谁用 & OpenClaw（PM AI） & Cursor Agent \\
格式 & 普通 Markdown & MDC（YAML + MD） \\
加载触发 & AI 自主判断 & \texttt{alwaysApply} 或 \texttt{globs} \\
记忆 & 有（memory 文件系统） & 无（每次对话全新） \\
通信 & 有（飞书通道） & 无 \\
代码修改 & 不直接改代码 & 读写代码、运行命令 \\
\bottomrule
\end{tabular}
\end{table}

\textbf{关键理解}：\texttt{skills/} 目录下的 SKILL.md 是给 \textbf{OpenClaw} 读的（让 PM 知道规范长什么样），\texttt{.cursor/rules/} 是给 \textbf{Cursor} 读的（让 Cursor 写代码时自动遵守规范）。两者内容有关联但职责不同。

% ---------- 第三章 ----------
\section{分工矩阵}

\begin{table}[h]
\centering
\small
\begin{tabular}{lcc}
\toprule
\textbf{职责} & \textbf{OpenClaw (PM)} & \textbf{Cursor Agent} \\
\midrule
需求生成与拆解   & \textcolor{success}{\checkmark} & \\
独立验收         & \textcolor{success}{\checkmark} & \\
文档维护与汇报   & \textcolor{success}{\checkmark} & \\
进度管理         & \textcolor{success}{\checkmark} & \\
UI 设计          & & \textcolor{success}{\checkmark} \\
前端开发         & & \textcolor{success}{\checkmark} \\
后端开发         & & \textcolor{success}{\checkmark} \\
环境搭建与运维   & & \textcolor{success}{\checkmark} \\
自验收           & & \textcolor{success}{\checkmark} \\
响应 PM 验收     & & \textcolor{success}{\checkmark} \\
\bottomrule
\end{tabular}
\end{table}

\subsection{自验收 vs 独立验收}

\textbf{Cursor 自验收}（每次开发完成后执行）：
\begin{itemize}[nosep]
  \item 运行 \texttt{npm run build} 确认前端无报错
  \item 运行后端启动确认服务正常
  \item 对照 testing.mdc 逐项自检
\end{itemize}

\textbf{PM 独立验收}（Cursor 返回结果后执行）：
\begin{itemize}[nosep]
  \item 检查产出物是否完整覆盖需求
  \item 站在用户角度验证功能
  \item 执行查证三问（目标达成？细节硬伤？换位质疑？）
  \item \textbf{不依赖} Cursor 的自测结论
\end{itemize}

% ---------- 第四章 ----------
\section{调用链路详解}

一次完整的任务执行流程：

\begin{enumerate}[nosep]
  \item \textbf{魏子政在飞书发消息}：如「加一个应用评分功能」
  \item \textbf{OpenClaw 接收}：加载 pm/SKILL.md，读取 memory 恢复上下文，读取项目文档
  \item \textbf{PM 执行}：记录任务到 memory，拆解需求为开发任务，写入 PROGRESS.md
  \item \textbf{调度 Cursor}：
\begin{lstlisting}
cursor agent --print --trust \
  --workspace projects/enterprise-app-store \
  "实现应用评分功能：1.数据库加 ratings 表
   2.后端评分CRUD接口 3.前端评分展示组件"
\end{lstlisting}
  \item \textbf{Cursor 执行}：自动加载 CLAUDE.md + rules/，编写代码，返回结果
  \item \textbf{PM 审计}：加载 tester/SKILL.md，检查产出物，执行查证三问
  \item \textbf{汇报}：飞书通知魏子政，更新 PROGRESS.md
\end{enumerate}

\subsection{CLI 调用参数说明}

\begin{lstlisting}
cursor agent --print --trust --workspace <项目路径> "任务描述"
# --print    : 输出完整响应（OpenClaw 可捕获）
# --trust    : 跳过人工确认（PM 调度场景）
# --workspace: 指定项目根目录
\end{lstlisting}

% ---------- 第五章 ----------
\section{OpenClaw 记忆系统}

Cursor \textbf{没有跨会话记忆}。每次对话都是全新的，只靠 CLAUDE.md 和 rules/ 获取上下文。

OpenClaw \textbf{有文件级记忆}：

\begin{table}[h]
\centering
\begin{tabular}{lp{8cm}l}
\toprule
\textbf{文件} & \textbf{用途} & \textbf{更新频率} \\
\midrule
SOUL.md & 身份、性格、行为准则 & 很少改 \\
USER.md & 用户信息、偏好 & 偶尔更新 \\
MEMORY.md & 长期记忆（项目经验、教训） & 每隔几天整理 \\
memory/YYYY-MM-DD.md & 当日任务日志 & 每天 \\
PROGRESS.md & 项目进度 & 每次任务后 \\
BUGS.md & Bug 记录 & 发现 Bug 时 \\
\bottomrule
\end{tabular}
\end{table}

这就是为什么 OpenClaw 能做 PM 而 Cursor 不能 -- Cursor 不知道昨天做了什么、上周改了什么需求、魏子政有什么偏好。OpenClaw 通过文件记忆「知道」这些。

% ---------- 第六章 ----------
\section{Skill/规则文件清单}

\subsection{OpenClaw 的 Skill}

\begin{table}[h]
\centering
\begin{tabular}{lll}
\toprule
\textbf{Skill} & \textbf{文件} & \textbf{用途} \\
\midrule
PM 产品经理 & \texttt{skills/pm/SKILL.md} & 需求拆解、任务分配、审计、查证 \\
UI 设计 & \texttt{skills/ui-design/SKILL.md} & 设计方案输出、视觉规范 \\
测试 & \texttt{skills/tester/SKILL.md} & 阶段门控验证、功能测试 \\
开发规范 & \texttt{skills/developer/SKILL.md} & 开发规范参考（PM 审计用） \\
\bottomrule
\end{tabular}
\end{table}

\subsection{Cursor 的规则文件}

\begin{table}[h]
\centering
\begin{tabular}{lll}
\toprule
\textbf{规则文件} & \textbf{触发条件} & \textbf{用途} \\
\midrule
\texttt{CLAUDE.md} & 每次必加载 & 项目总入口，指向 rules/ \\
\texttt{project-overview.mdc} & alwaysApply & 技术栈锁定、禁止事项 \\
\texttt{backend-rules.mdc} & backend/ & async/await、Pydantic、错误格式 \\
\texttt{frontend-rules.mdc} & frontend/ & \texttt{<script setup>}、TS、Pinia \\
\texttt{ui-design.mdc} & frontend/ & 配色、间距、组件规范 \\
\texttt{api-conventions.mdc} & backend/ & RESTful 命名、分页、统一响应 \\
\texttt{testing.mdc} & alwaysApply & 自测要求 \\
\texttt{database.mdc} & 数据库相关 & 迁移、软删除、审计字段 \\
\bottomrule
\end{tabular}
\end{table}

% ============================================================
%  第二部分：技术选型
% ============================================================
\part{技术选型}

% ---------- 第七章 ----------
\section{锁定技术栈总览}

\begin{table}[h]
\centering
\begin{tabular}{llp{5.5cm}}
\toprule
\textbf{层} & \textbf{技术} & \textbf{用途} \\
\midrule
前端 & Vue 3 + Vite + Element Plus + Pinia + TypeScript & 管理后台 + 开发者门户 \\
后端 & FastAPI + SQLAlchemy (async) + Pydantic & REST API 服务 \\
数据库 & PostgreSQL 16 & 元数据 + JSONB + 审计日志 \\
对象存储 & MinIO (S3 兼容) & 应用包存储 \\
搜索 & OpenSearch 2.x & 全文检索 + 向量语义搜索 \\
认证 & JWT (python-jose + passlib) & 用户认证 + RBAC \\
容器化 & Docker + Docker Compose & 本地开发 + 生产部署 \\
迁移 & Alembic & 数据库 schema 版本管理 \\
\bottomrule
\end{tabular}
\end{table}

\textbf{硬约束：前端不用 Next.js/React，后端不用 Spring/Java。技术栈一经锁定，OpenClaw 和 Cursor 均不得自行变更。}

% ---------- 第八章 ----------
\section{逐项选型原因}

\subsection{前端：Vue 3 + Element Plus}

\textbf{为什么选 Vue 3？}

\begin{enumerate}[nosep]
  \item 学习曲线低：模板语法直观，实习生上手快，比 React hooks 概念负担轻
  \item Element Plus 生态成熟：企业级中后台组件库，表格/表单/弹窗/分页开箱即用，减少 50\%+ UI 工作量
  \item 国内团队主流：中文文档完善，社区活跃，招人容易
  \item TypeScript 一等支持：\texttt{<script setup>} + TS 组合式 API
  \item Vite 构建极快：HMR 毫秒级
\end{enumerate}

\textbf{为什么不用 Next.js/React？} SSR/SSG 对内部管理系统价值不大；React 需额外状态管理方案；团队偏 Python，Vue 模板语法更易跨角色理解。

\textbf{Element Plus vs Ant Design Vue vs Arco Design}：Element Plus 组件最全、文档最好、社区最大，选定。

\subsection{后端：FastAPI}

\textbf{为什么选 FastAPI？}

\begin{enumerate}[nosep]
  \item 异步原生：async/await 一等支持，天然适合 IO 密集型（文件上传、搜索调用、扫描集成）
  \item 自动文档：Pydantic 模型自动生成 Swagger/OpenAPI 文档
  \item 类型安全：请求/响应自动校验，参数错误自动返回 422
  \item Python 生态：与安全扫描工具（Semgrep/ClamAV）、LLM API 天然互通
  \item 开发效率高：代码量少，原型快
\end{enumerate}

\textbf{为什么不用 Spring Boot？} 启动慢、代码量大，对中小团队效率偏低；安全扫描工具 Python 集成更方便。

\subsection{数据库：PostgreSQL 16}

\textbf{为什么选 PostgreSQL？}

\begin{enumerate}[nosep]
  \item 事务完整：ACID 保证，适合应用元数据、权限、审计
  \item JSONB 类型：扫描结果是嵌套 JSON，PG JSONB 支持查询和索引
  \item 扩展性强：pgvector（向量搜索）、全文检索、范围类型
  \item 运维成熟：备份、复制、监控工具链完善
\end{enumerate}

\textbf{为什么不用 MongoDB？} 去掉 MongoDB 减少组件；扫描结果用 PG JSONB 足够。\\
\textbf{为什么不用 MySQL？} JSONB 支持不如 PG 灵活，复杂查询能力偏弱。

\subsection{搜索引擎：OpenSearch 2.x}

\textbf{为什么选 OpenSearch？}

\begin{enumerate}[nosep]
  \item 全文 + 向量一体化：一个集群同时支持关键词搜索和向量语义搜索
  \item 权限过滤：DSL 支持 per-document 权限查询
  \item 中文分词：IK Analysis 插件
  \item Apache 2.0 许可：无 SSPL 限制（Elasticsearch 7.11+ 改为 SSPL）
\end{enumerate}

\subsection{对象存储：MinIO}

\textbf{为什么选 MinIO？}

\begin{enumerate}[nosep]
  \item S3 API 兼容：随时可切换到 AWS S3/阿里 OSS
  \item 私有化部署：数据不出域
  \item 预签名 URL：前端直接下载，不经后端转发
  \item Docker 部署简单：一个容器搞定
\end{enumerate}

\subsection{认证授权：JWT + RBAC}

\textbf{为什么选 JWT 而不是 Keycloak？}

\begin{enumerate}[nosep]
  \item 轻量：Keycloak 需要 JVM + 数据库 + 独立服务，JWT 只需两个 Python 库
  \item 够用：初期只需 4 角色（admin/reviewer/developer/viewer）
  \item 无状态：Token 自包含，不需要每次查库
\end{enumerate}

\textbf{后续扩展}：需要 LDAP/SSO/多租户时引入 Keycloak，JWT 签发逻辑从本地改为对接 Keycloak。

\subsection{容器化：Docker + Docker Compose}

\textbf{为什么选 Docker Compose 而不是 K8s？}

\begin{enumerate}[nosep]
  \item MVP 不需要 K8s 编排/调度/自愈能力
  \item 开发体验好：\texttt{docker compose up} 一键启动所有依赖
  \item 运维简单：5--10 人团队不需要专门 K8s 运维
\end{enumerate}

\textbf{后续扩展}：需要多实例/灰度发布时迁移 K8s，Dockerfile 不需要改，只需加 Helm Chart。

% ---------- 第九章 ----------
\section{技术栈演进记录}

\begin{table}[h]
\centering
\begin{tabular}{llp{5cm}}
\toprule
\textbf{原组件} & \textbf{替代方案} & \textbf{精简原因} \\
\midrule
Next.js & Vue 3 + Vite & 内部管理不需要 SSR \\
MongoDB & PostgreSQL JSONB & 减少组件，JSONB 足够 \\
Keycloak & JWT + RBAC 表 & 初期不需要完整 IdP \\
Temporal & 数据库状态机 & 审批流程简单，不需要工作流引擎 \\
Kafka & Python asyncio & 初期异步量不大 \\
K8s + Helm & Docker Compose & MVP 单机部署足够 \\
Prometheus + Grafana & 延后实现 & MVP 优先功能，监控后续补 \\
\bottomrule
\end{tabular}
\caption{v1.0 $\to$ v2.0 精简记录}
\end{table}

\textbf{原则}：「够用即可、按需扩展」。后续需要时再逐个引入。

% ============================================================
%  第三部分：文档体系
% ============================================================
\part{文档体系与生命周期}

% ---------- 第十章 ----------
\section{文档目录结构}

\begin{lstlisting}
enterprise-app-store/
  reference/                          # 参考文档（调研、规范）
    DOC-NAMING-CONVENTION.tex/pdf    # 文档命名规范
    REF-001-skill-store-arch-survey  # 开源 Skill 商店架构调研
    ...                              # 后续新增 REF-xxx
  demands/                            # 需求变更文档
    DEMAND-001-role-redefine         # 角色分工重新定义
    ...                              # 后续新增 DEMAND-xxx
  AGENT-ARCHITECTURE.md               # 智能体搭配架构
  tech-selection.md                   # 技术选型说明书
  TEAM-MANUAL.md                      # 团队协作说明书
  DOC-LIFECYCLE.md                    # 文档生命周期管理
  CHANGELOG.md                        # 变更日志
  PROGRESS.md                         # 进度追踪表
  BUGS.md                             # Bug 记录
  CLAUDE.md                           # Cursor 总入口
  .cursor/rules/*.mdc                 # Cursor 编码规则
  skills/pm|tester|ui-design|developer/SKILL.md
  database/SCHEMA.md                  # 数据库设计
\end{lstlisting}

\subsection{文档分类与维护责任}

\begin{table}[h]
\centering
\small
\begin{tabular}{llp{4.5cm}}
\toprule
\textbf{类别} & \textbf{维护者} & \textbf{文档举例} \\
\midrule
产品与规范 & OpenClaw & AGENT-ARCHITECTURE.md, tech-selection.md, TEAM-MANUAL.md \\
参考文档 (reference/) & OpenClaw & REF-xxx, DOC-NAMING-CONVENTION \\
需求文档 (demands/) & OpenClaw & DEMAND-xxx \\
设计与规范 & OpenClaw 设计 & SCHEMA.md, skills/*/SKILL.md \\
Cursor 规则 & OpenClaw & CLAUDE.md, .cursor/rules/*.mdc \\
开发文档 & Cursor 输出 & BUGS.md \\
\bottomrule
\end{tabular}
\end{table}

% ---------- 第十一章 ----------
\section{tex-pdf 文档规范}

\subsection{格式要求}

\begin{itemize}[nosep]
  \item \textbf{源文件}：\texttt{.tex} 格式，提交 Git 版本控制
  \item \textbf{交付物}：\texttt{.pdf} 格式，用于审阅和分发
  \item \textbf{编译工具}：Tectonic（\texttt{\textasciitilde/.local/tectonic/tectonic}）
  \item \textbf{管理者}：PM AI（万能产品小助手）
\end{itemize}

\subsection{文件命名规范}

格式：\texttt{\{前缀\}-\{三位编号\}-\{kebab-case 简述\}.tex/pdf}

\begin{table}[h]
\centering
\begin{tabular}{lll}
\toprule
\textbf{前缀} & \textbf{含义} & \textbf{目录} \\
\midrule
REF & 参考文档 & reference/ \\
DEMAND & 需求变更 & demands/ \\
SPEC & 技术规格 & reference/ \\
GUIDE & 操作指南 & reference/ \\
\bottomrule
\end{tabular}
\end{table}

\subsection{编译命令}

\begin{lstlisting}
cd reference/ && ~/.local/tectonic/tectonic REF-001-xxx.tex
cd demands/ && ~/.local/tectonic/tectonic DEMAND-001-xxx.tex
\end{lstlisting}

\subsection{已知限制}

\begin{enumerate}[nosep]
  \item 避免 fontawesome5 命令（\texttt{\textbackslash faCalendar} 不可用）
  \item 避免 Unicode 箭头（用 \texttt{\textbackslash to} 或 \texttt{-{}->{}} 替代）
  \item 避免 listings 的 \texttt{language} 参数（部分语言定义缺失）
  \item Overfull hbox 少量警告可接受
\end{enumerate}

% ---------- 第十二章 ----------
\section{文档版本管理}

\subsection{版本号规则}

语义化版本：\texttt{vX.Y.Z}

\begin{table}[h]
\centering
\begin{tabular}{cll}
\toprule
\textbf{位} & \textbf{含义} & \textbf{示例} \\
\midrule
X（主版本） & 架构/技术栈重大变更 & v1.0 $\to$ v2.0 \\
Y（次版本） & 新增模块/功能 & v2.0 $\to$ v2.1 \\
Z（补丁） & 文档修正、措辞优化 & v2.1 $\to$ v2.1.1 \\
\bottomrule
\end{tabular}
\end{table}

\subsection{版本阶段定义}

\begin{table}[h]
\centering
\begin{tabular}{clp{6cm}}
\toprule
\textbf{阶段} & \textbf{名称} & \textbf{内容范围} \\
\midrule
V0 & 先备阶段 & 智能体分工、商店整体框架、文档体系、技术选型 \\
V0.x & 先备迭代 & 分工调整、架构调研、参考方案、规范补充 \\
V1 & 正式需求 & 第一批正式业务需求接入 \\
V1.x & 功能迭代 & 功能增强、Bug 修复、性能优化 \\
V2+ & 演进阶段 & 架构升级、生态扩展、商业运营 \\
\bottomrule
\end{tabular}
\end{table}

\textbf{当前阶段}：V0（先备阶段）。本总纲是 V0 的核心交付物。

\subsection{变更日志}

项目根目录维护统一的 \texttt{CHANGELOG.md}，记录所有文档和代码的变更。每个重要文档头部标注版本号、日期、变更摘要。

% ---------- 第十三章 ----------
\section{需求新增处理流程}

\subsection{需求来源}

\begin{table}[h]
\centering
\begin{tabular}{ll}
\toprule
\textbf{来源} & \textbf{示例} \\
\midrule
魏子政直接提 & 「加一个应用评分功能」 \\
阶段审计发现 & 「缺少数据导出功能」 \\
Bug 修复衍生 & 「修复后发现需要加校验」 \\
\bottomrule
\end{tabular}
\end{table}

\subsection{5 步处理流程}

\begin{enumerate}
  \item \textbf{需求记录}：写入 memory 文件，包含需求描述、提出者、优先级
  \item \textbf{影响评估}（OpenClaw 执行）：
    \begin{itemize}[nosep]
      \item 技术影响：改哪些模块？新增哪些表/接口/页面？
      \item 文档影响：哪些文档需要更新？
      \item 进度影响：对当前里程碑有无影响？
    \end{itemize}
  \item \textbf{方案设计}（复杂需求）：更新 SCHEMA.md、补充 SKILL.md、涉及技术栈变更须魏子政审批
  \item \textbf{任务拆解与分配}：拆解为可独立完成的开发任务，写入 PROGRESS.md
  \item \textbf{开发--测试--审计--交付}：Cursor 执行 $\to$ OpenClaw 审计 $\to$ 文档同步 $\to$ 汇报
\end{enumerate}

\subsection{优先级定义}

\begin{table}[h]
\centering
\begin{tabular}{clll}
\toprule
\textbf{级别} & \textbf{标签} & \textbf{含义} & \textbf{处理时效} \\
\midrule
P0 & \textcolor{danger}{阻塞} & 阻塞当前阶段推进 & 立即 \\
P1 & \textcolor{warning}{严重} & 核心功能缺陷 & 当天 \\
P2 & 一般 & 重要但不紧急 & 排入当前/下阶段 \\
P3 & 轻微 & 锦上添花 & 有空再做 \\
\bottomrule
\end{tabular}
\end{table}

\subsection{需求变更控制}

\textbf{小变更}（不影响已有模块）：PM 直接处理，事后说明。\\
\textbf{中变更}（影响已有模块但不涉及架构）：PM 评估 $\to$ 简要说明 $\to$ 执行。\\
\textbf{大变更}（涉及架构/技术栈/数据库结构）：PM 评估 $\to$ 详细方案 $\to$ \textbf{魏子政审批} $\to$ 执行。

% ============================================================
%  第四部分：自闭环机制
% ============================================================
\part{自闭环机制}

% ---------- 第十四章 ----------
\section{自闭环定义}

自闭环 = 从需求提出到交付完成，\textbf{所有环节在同一体系内完成}，不需要外部干预。

\begin{lstlisting}
魏子政提需求
  -> OpenClaw 接收 -> 记录 -> 评估
  -> 方案设计 -> 文档更新
  -> Cursor 开发 -> 代码产出
  -> OpenClaw 测试 -> 查证 -> 审计
  -> 文档同步更新
  -> 交付输出 -> 反馈 -> 循环
\end{lstlisting}

% ---------- 第十五章 ----------
\section{五个闭环}

\subsection{闭环 1：需求闭环}

\begin{itemize}[nosep]
  \item \textbf{输入}：魏子政提需求
  \item \textbf{过程}：记录 $\to$ 评估 $\to$ 设计 $\to$ 拆解 $\to$ 分配
  \item \textbf{输出}：PROGRESS.md 中的任务条目（附验收标准）
  \item \textbf{闭合标志}：PROGRESS.md 出现新任务，状态为「待开始」
\end{itemize}

\subsection{闭环 2：开发闭环}

\begin{itemize}[nosep]
  \item \textbf{输入}：Cursor 收到任务
  \item \textbf{过程}：读 rules/ $\to$ 写代码 $\to$ 自测 $\to$ 提交
  \item \textbf{输出}：可运行的代码
  \item \textbf{闭合标志}：\texttt{npm run build} 通过 + 后端启动正常
\end{itemize}

\subsection{闭环 3：测试闭环}

\begin{itemize}[nosep]
  \item \textbf{输入}：Cursor 产出代码
  \item \textbf{过程}：OpenClaw 执行测试验证（接口 $\to$ 功能 $\to$ 流程 $\to$ UI）
  \item \textbf{输出}：测试报告
  \item \textbf{闭合标志}：tester/SKILL.md 验证项全部通过
\end{itemize}

\subsection{闭环 4：文档闭环}

\begin{itemize}[nosep]
  \item \textbf{输入}：代码/需求/规范变更
  \item \textbf{过程}：更新受影响文档 $\to$ 生成新 PDF $\to$ 记录 CHANGELOG
  \item \textbf{输出}：所有文档与代码保持一致
  \item \textbf{闭合标志}：CHANGELOG.md 有新记录，文档版本号已更新
\end{itemize}

\subsection{闭环 5：审计闭环}

\begin{itemize}[nosep]
  \item \textbf{输入}：测试通过的开发任务
  \item \textbf{过程}：OpenClaw 执行查证三问（目标达成？细节硬伤？换位质疑？）
  \item \textbf{输出}：审计结论（通过/需修正）
  \item \textbf{闭合标志}：查证三问全部通过，PROGRESS.md 任务状态变为「已完成」
\end{itemize}

% ---------- 第十六章 ----------
\section{闭环监控标准}

OpenClaw 在每次任务完成后，检查所有 5 个闭环是否闭合：

\begin{table}[h]
\centering
\begin{tabular}{lp{8cm}}
\toprule
\textbf{检查项} & \textbf{闭合标准} \\
\midrule
需求闭环 & PROGRESS.md 有任务记录 \\
开发闭环 & Cursor 返回成功结果 \\
测试闭环 & tester/SKILL.md 验证项全部通过 \\
文档闭环 & 受影响文档已更新 + CHANGELOG.md 有记录 \\
审计闭环 & 查证三问通过 \\
\bottomrule
\end{tabular}
\end{table}

\textbf{任何闭环未闭合，任务不能标记为「已完成」。}

% ============================================================
%  第五部分：当前状态与下一步
% ============================================================
\part{当前状态与下一步}

% ---------- 第十七章 ----------
\section{V0 阶段已完成项}

\begin{enumerate}[nosep]
  \item \textbf{智能体分工定义}：OpenClaw = 需求生成+拆解+独立验收；Cursor = 前后端+UI+环境+运维+测试
  \item \textbf{技术选型 v2.0 锁定}：Vue 3 + FastAPI + PG + MinIO + OpenSearch + JWT + Docker Compose
  \item \textbf{Skill/规则体系搭建}：
    \begin{itemize}[nosep]
      \item OpenClaw：4 个 SKILL.md（pm、tester、ui-design、developer）
      \item Cursor：1 个 CLAUDE.md + 7 个 .mdc 规则文件
    \end{itemize}
  \item \textbf{文档体系搭建}：
    \begin{itemize}[nosep]
      \item 产品规范：AGENT-ARCHITECTURE.md、tech-selection.md、TEAM-MANUAL.md、DOC-LIFECYCLE.md
      \item tex-pdf 规范：DOC-NAMING-CONVENTION（命名规范）、REF-001（架构调研）
      \item 需求文档：DEMAND-001（角色分工重新定义）
    \end{itemize}
  \item \textbf{自闭环机制设计}：5 个闭环 + 监控标准
  \item \textbf{数据库设计}：database/SCHEMA.md（ER 图 + 数据字典）
  \item \textbf{开源方案调研}：Flow Nexus App Store（MCP 架构）+ Membrane Marketplacer（中间件代理）
\end{enumerate}

% ---------- 第十八章 ----------
\section{V0 阶段待完成项}

\begin{enumerate}[nosep]
  \item 强化 \texttt{ui-design.mdc}，补充常用页面模式（列表页、详情页、表单页、审批页）
  \item 从 cursor.directory 精选 a11y/响应式相关规则片段
  \item 评估 Figma MCP 集成可行性（如魏子政提供设计稿）
  \item 确认数据库 schema 与技术选型的最终一致性
\end{enumerate}

% ---------- 第十九章 ----------
\section{V1 阶段规划}

V1 阶段开始接入正式业务需求，预计里程碑：

\begin{table}[h]
\centering
\begin{tabular}{lll}
\toprule
\textbf{阶段} & \textbf{核心交付物} & \textbf{预计周期} \\
\midrule
环境搭建 & Docker Compose + 登录页 + /health & 1 周 \\
核心功能 & 应用 CRUD + 文件上传 + 用户认证 & 1 周 \\
搜索 + 审批 & OpenSearch + 审批流程 & 1 周 \\
安全扫描 & Semgrep + ClamAV + LLM 集成 & 1 周 \\
打磨 + 部署 & UI 优化 + Docker 生产部署 & 2 周 \\
\bottomrule
\end{tabular}
\end{table}

% ---------- 附录 ----------
\appendix
\section{附录 A：Git 分支策略}

\begin{lstlisting}
main          # 稳定版本，里程碑合并
dev           # 开发分支，日常提交
  feat/xxx    # 新功能分支
  fix/xxx     # Bug 修复分支
\end{lstlisting}

\section{附录 B：提交规范}

\begin{lstlisting}
feat: 新增应用评分功能
fix: 修复登录页 token 过期未跳转
docs: 更新技术选型文档 v2.0
refactor: 重构上传模块为异步
\end{lstlisting}

\section{附录 C：新成员上手阅读顺序}

\begin{enumerate}[nosep]
  \item TEAM-MANUAL.md -- 了解怎么协作
  \item AGENT-ARCHITECTURE.md -- 了解智能体怎么配合
  \item tech-selection.md -- 了解技术栈和选型原因
  \item database/SCHEMA.md -- 了解数据库设计
  \item PROGRESS.md -- 了解当前进度
\end{enumerate}

\section{附录 D：文档索引}

\begin{table}[h]
\centering
\small
\begin{tabular}{llp{6cm}}
\toprule
\textbf{文档} & \textbf{类型} & \textbf{说明} \\
\midrule
本文件（项目总纲 PDF） & tex-pdf & V0 核心交付物，总纲 \\
AGENT-ARCHITECTURE.md & Markdown & 智能体搭配详细说明 \\
tech-selection.md & Markdown & 技术选型 v2.0 \\
TEAM-MANUAL.md & Markdown & 团队协作说明书 \\
DOC-LIFECYCLE.md & Markdown & 文档生命周期管理 \\
CHANGELOG.md & Markdown & 变更日志 \\
PROGRESS.md & Markdown & 进度追踪表 \\
CLAUDE.md & Markdown & Cursor 总入口 \\
database/SCHEMA.md & Markdown & 数据库设计 \\
BUGS.md & Markdown & Bug 记录 \\
reference/DOC-NAMING-CONVENTION & tex-pdf & 文档命名规范 \\
reference/REF-001-skill-store-arch-survey & tex-pdf & 开源 Skill 商店架构调研 \\
demands/DEMAND-001-role-redefine & tex-pdf & 角色分工重新定义 \\
\bottomrule
\end{tabular}
\end{table}

\end{document}
